\documentclass[11pt,a4paper]{article}
\usepackage{fontspec,amsmath,amssymb,amsthm,microtype}
\setmainfont{lmroman10-regular.otf}[BoldFont=lmroman10-bold.otf,ItalicFont=lmroman10-italic.otf,BoldItalicFont=lmroman10-bolditalic.otf]
\setmonofont{lmmono10-regular.otf}[Scale=MatchLowercase]
\usepackage[margin=25mm]{geometry}
\usepackage[unicode,hidelinks]{hyperref}
\setlength{\emergencystretch}{3em}
\setlength{\parskip}{0.4em}
\setlength{\parindent}{0pt}
\allowdisplaybreaks[2]
\raggedbottom
\title{\texorpdfstring{The genuine frozen-\allowbreak{}linear contraction in Proposition 3}{The genuine frozen-linear contraction in Proposition 3}}
\author{LangevinLean mathematical exposition}
\date{}
\begin{document}\maketitle
\noindent Source module: \texttt{Kinetic\allowbreak{}Langevin/\allowbreak{}Second\allowbreak{}Order\allowbreak{}Linear.\allowbreak{}lean}. The accompanying HTML pairs this mathematical exposition with the exact Lean statements and records the scope of the formalization.


The friction can be much larger than the square root of the smoothness constant. A crude Euler remainder proportional to \ensuremath{\gamma}\ensuremath{{}^{2}} would lose the paper’s step-size regime. The identities below instead use the force’s cocoercivity to retain a second-order budget proportional to M, and prove the actual frozen-linear map contracts by 1\ensuremath{-}mt/\allowbreak{}(2\ensuremath{\gamma}) under 5\ensuremath{\gamma}Mt\ensuremath{\le}m.


\section{\texorpdfstring{Separate fast damping from the Euler remainder}{Separate fast damping from the Euler remainder}}


For the kinetic difference drift A=(\ensuremath{-}\ensuremath{\gamma}v\ensuremath{-}F,v), its transformed squared norm satisfies the exact identity E(A)=\ensuremath{-}\ensuremath{\gamma}E\ensuremath{\prime}/\allowbreak{}2+2\ensuremath{\Vert}F\ensuremath{\Vert}\ensuremath{{}^{2}}\ensuremath{-}\ensuremath{\gamma}\ensuremath{{}^{2}}\ensuremath{\langle}x,F\ensuremath{\rangle}. The actual Euler phase (v,x)+sA has energy E+sE\ensuremath{\prime}+s\ensuremath{{}^{2}}E(A). Lean expands the two transformed coordinates, uses the real inner-product identities, and checks both equalities by ring normalization.



{\raggedright\noindent Lean declarations: \texttt{kinetic\_\allowbreak{}drift\_\allowbreak{}energy\_\allowbreak{}identity}, \texttt{kinetic\_\allowbreak{}euler\_\allowbreak{}energy\_\allowbreak{}identity}.\par}

\section{\texorpdfstring{Use cocoercivity before bounding the remainder}{Use cocoercivity before bounding the remainder}}


The force satisfies \ensuremath{\Vert}F\ensuremath{\Vert}\ensuremath{{}^{2}}\ensuremath{\le}M\ensuremath{\langle}x,F\ensuremath{\rangle} and \ensuremath{\langle}x,F\ensuremath{\rangle}\ensuremath{\le}M\ensuremath{\Vert}x\ensuremath{\Vert}\ensuremath{{}^{2}}. Since \ensuremath{\gamma}\ensuremath{{}^{2}}\ensuremath{\ge}m+M, 2M\ensuremath{-}\ensuremath{\gamma}\ensuremath{{}^{2}}\ensuremath{\le}M\ensuremath{-}m. Also \ensuremath{\gamma}\ensuremath{{}^{2}}\ensuremath{\Vert}x\ensuremath{\Vert}\ensuremath{{}^{2}}\ensuremath{\le}2E. Combining these facts gives E(A)\ensuremath{\le}\ensuremath{-}\ensuremath{\gamma}E\ensuremath{\prime}/\allowbreak{}2+2M(M\ensuremath{-}m)E/\allowbreak{}\ensuremath{\gamma}\ensuremath{{}^{2}}. In the Euler expansion, the derivative coefficient becomes s(1\ensuremath{-}\ensuremath{\gamma}s/\allowbreak{}2), which is nonnegative for \ensuremath{\gamma}s\ensuremath{\le}1. Actual force dissipativity E\ensuremath{\prime}\ensuremath{\le}\ensuremath{-}2mE/\allowbreak{}\ensuremath{\gamma} then yields transformed norm at most [1\ensuremath{-}ms/\allowbreak{}\ensuremath{\gamma}+(m/\allowbreak{}2+M(M\ensuremath{-}m)/\allowbreak{}\ensuremath{\gamma}\ensuremath{{}^{2}})s\ensuremath{{}^{2}}] times the initial norm. The final square comparison retains an additional nonnegative square instead of estimating it away.



{\raggedright\noindent Lean declarations: \texttt{kinetic\_\allowbreak{}drift\_\allowbreak{}energy\_\allowbreak{}bound}, \texttt{kinetic\_\allowbreak{}euler\_\allowbreak{}norm\_\allowbreak{}bound}.\par}

\section{\texorpdfstring{Bound the true frozen-\allowbreak{}Hessian correction}{Bound the true frozen-Hessian correction}}


The remaining phase correction is (\ensuremath{-}\ensuremath{\varphi}\ensuremath{{}_{2}}Hv,\ensuremath{-}\ensuremath{\psi}\ensuremath{{}_{2}}F\ensuremath{-}\ensuremath{\varphi}\ensuremath{{}_{3}}Hv). The exact coefficient identity \ensuremath{\varphi}\ensuremath{{}_{2}}+\ensuremath{\gamma}\ensuremath{\varphi}\ensuremath{{}_{3}}=\ensuremath{\psi}\ensuremath{{}_{2}} transforms it into (\ensuremath{-}\ensuremath{\psi}\ensuremath{{}_{2}}(Hv+\ensuremath{\gamma}F),\ensuremath{\varphi}\ensuremath{{}_{2}}Hv). The bounds \ensuremath{\psi}\ensuremath{{}_{2}},\ensuremath{\varphi}\ensuremath{{}_{2}}\ensuremath{\le}t\ensuremath{{}^{2}}/\allowbreak{}2, \ensuremath{\Vert}H\ensuremath{\Vert}\ensuremath{\le}M, \ensuremath{\Vert}v\ensuremath{\Vert}\ensuremath{\le}N, and \ensuremath{\gamma}\ensuremath{\Vert}x\ensuremath{\Vert}\ensuremath{\le}3N/\allowbreak{}2 give component envelopes 5Mt\ensuremath{{}^{2}}N/\allowbreak{}4 and Mt\ensuremath{{}^{2}}N/\allowbreak{}2. Their squared sum is below (3Mt\ensuremath{{}^{2}}N/\allowbreak{}2)\ensuremath{{}^{2}}. Thus the actual correction costs at most 3Mt\ensuremath{{}^{2}}N/\allowbreak{}2, with no symmetry or common Hessian eigenbasis assumed.



{\raggedright\noindent Lean declarations: \texttt{second\allowbreak{}Order\_\allowbreak{}frozen\_\allowbreak{}correction\_\allowbreak{}bound}.\par}

\section{\texorpdfstring{Check the scalar coefficient budget with exact \ensuremath{\psi}\ensuremath{{}_{1}}}{Check the scalar coefficient budget with exact ψ₁}}


The friction balance \ensuremath{\gamma}\ensuremath{\psi}\ensuremath{{}_{2}}=t\ensuremath{-}\ensuremath{\psi}\ensuremath{{}_{1}} and \ensuremath{\psi}\ensuremath{{}_{2}}\ensuremath{\le}t\ensuremath{{}^{2}}/\allowbreak{}2 imply \ensuremath{\psi}\ensuremath{{}_{1}}\ensuremath{\ge}t\ensuremath{-}\ensuremath{\gamma}t\ensuremath{{}^{2}}/\allowbreak{}2, while 0\ensuremath{\le}\ensuremath{\psi}\ensuremath{{}_{1}}\ensuremath{\le}t. Moreover m+M(M\ensuremath{-}m)/\allowbreak{}\ensuremath{\gamma}\ensuremath{{}^{2}}\ensuremath{\le}M. These inequalities bound the Euler coefficient plus the frozen correction by 1\ensuremath{-}mt/\allowbreak{}\ensuremath{\gamma}+5Mt\ensuremath{{}^{2}}/\allowbreak{}2. The source step restriction 5\ensuremath{\gamma}Mt\ensuremath{\le}m makes the last term at most mt/\allowbreak{}(2\ensuremath{\gamma}), leaving 1\ensuremath{-}mt/\allowbreak{}(2\ensuremath{\gamma}). All divisions are justified by positive friction.



{\raggedright\noindent Lean declarations: \texttt{second\allowbreak{}Order\_\allowbreak{}linear\_\allowbreak{}scalar\_\allowbreak{}bound}.\par}

\section{\texorpdfstring{Identify and contract the actual frozen-\allowbreak{}linear map}{Identify and contract the actual frozen-linear map}}


frozenSecondOrderStep is the explicit phase with \ensuremath{\psi}\ensuremath{{}_{0}}v\ensuremath{-}\ensuremath{\psi}\ensuremath{{}_{1}}F\ensuremath{-}\ensuremath{\varphi}\ensuremath{{}_{2}}Hv in velocity and x+\ensuremath{\psi}\ensuremath{{}_{1}}v\ensuremath{-}\ensuremath{\psi}\ensuremath{{}_{2}}F\ensuremath{-}\ensuremath{\varphi}\ensuremath{{}_{3}}Hv in position. The identity \ensuremath{\psi}\ensuremath{{}_{0}}=1\ensuremath{-}\ensuremath{\gamma}\ensuremath{\psi}\ensuremath{{}_{1}} writes this actual map as the Euler phase plus the correction above. Strong-smooth interpolation supplies both cocoercivity and force dissipativity. The transformed triangle inequality and the checked scalar budget therefore prove the source contraction factor for the actual frozen-linear part of KLMC2.



{\raggedright\noindent Lean declarations: \texttt{frozen\allowbreak{}Second\allowbreak{}Order\allowbreak{}Step}, \texttt{frozen\allowbreak{}Second\allowbreak{}Order\allowbreak{}Step\_\allowbreak{}decomposition}, \texttt{frozen\allowbreak{}Second\allowbreak{}Order\allowbreak{}Step\_\allowbreak{}contract}.\par}

\end{document}
